\chapter{TỔNG QUAN VỀ ĐỀ TÀI}
\label{chap:tongquan}

\section{Bối cảnh nghiên cứu và tính cấp thiết của đề tài}
\label{sec:boicanh}

\subsection{Quan sát ban đầu từ dữ liệu và động lực chọn đề tài}
\label{subsec:thuctrang}

Động lực trực tiếp của đề tài xuất phát từ việc khảo sát bộ dữ liệu H\&M Personalized Fashion Recommendations \parencite{hm2022kaggle}: tệp giao dịch gốc ghi nhận 31.788.324 dòng mua hàng của 1.362.281 khách hàng trên 104.547 sản phẩm thời trang trong hai năm (mục~\ref{subsec:hm_description}). Hai đặc điểm của dữ liệu này đặt ra yêu cầu cho mô hình gợi ý. Thứ nhất, ma trận tương tác \emph{cực thưa}: ngay cả sau khi lọc $k$-core, mật độ chỉ khoảng 0,03\% trên toàn bộ dữ liệu (Bảng~\ref{tab:kcore_hm_full}) -- mỗi khách chỉ mua một phần rất nhỏ của danh mục, và lệch mạnh về một nhóm nhỏ sản phẩm phổ biến. Thứ hai, nhu cầu thời trang \emph{thay đổi theo mùa và theo xu hướng}: sản phẩm bán chạy trong vài tuần gần đây chưa chắc là sản phẩm được mua nhiều nhất trong hai năm. Một mô hình chỉ học từ mã khách hàng và mã sản phẩm không dùng được những thông tin này -- xu hướng bán gần đây của sản phẩm, sản phẩm thuộc loại gì, khách hàng thuộc nhóm nào -- trong khi bộ dữ liệu có sẵn thuộc tính và mô tả văn bản của từng sản phẩm, thông tin khách hàng và doanh số theo thời gian.

Câu hỏi của đề tài vì vậy là: trên dữ liệu giao dịch thời trang thực tế -- cực thưa, lệch về sản phẩm phổ biến và thay đổi theo mùa -- một mô hình lai kết hợp nhân tử hóa ma trận với mạng nơ-ron sâu, được bổ sung đặc trưng sản phẩm, khách hàng và thời gian, có cho chất lượng gợi ý tốt hơn các mô hình đối chứng được tinh chỉnh đầy đủ và tốt hơn chính các thành phần của nó hay không. Đề tài không dựa vào kết quả công bố của các công trình tham khảo (dùng dữ liệu phim ảnh, mạng xã hội hình ảnh với giao thức khác), mà xây dựng và đánh giá mô hình trực tiếp trên dữ liệu H\&M với cách chia dữ liệu và giao thức đánh giá riêng.

Ở góc độ ứng dụng, cá nhân hóa gợi ý sản phẩm trong thương mại điện tử giúp giảm chi phí tìm kiếm của khách hàng và hỗ trợ doanh nghiệp dự báo nhu cầu \parencite{schafer2001commerce, linden2003amazon}. Với ngành thời trang, nơi nhu cầu thay đổi theo mùa, gợi ý còn cần phản ánh xu hướng theo thời gian chứ không chỉ sở thích lâu dài của khách.

\subsection{Vì sao nhân tử hóa ma trận tuyến tính chưa đủ cho dữ liệu này}
\label{subsec:hanche_cf}

Cách tiếp cận đơn giản nhất cho bài toán trên là Nhân tử hóa ma trận (Matrix Factorization -- MF): mỗi User $u$ và Item $i$ được gán một vector ẩn $p_u, q_i \in \mathbb{R}^d$, và mức độ tương tác được dự đoán bằng tích vô hướng của hai vector \parencite{koren2009matrix}:
\begin{equation}
    \hat{y}_{ui} = p_u^T q_i = \sum_{f=1}^{d} p_{uf} q_{if}
    \label{eq:mf_intro}
\end{equation}
Công thức này có hai giới hạn trên dữ liệu H\&M. Một là giới hạn về \emph{biểu diễn}: tích vô hướng cộng dồn các chiều đặc trưng theo cùng một trọng số cố định và chỉ mô hình hóa được tương tác tuyến tính -- He và cộng sự \parencite{he2017ncf} lấy đây làm động cơ để bổ sung một thành phần phi tuyến. Hai là giới hạn về \emph{thông tin}: vector $q_i$ chỉ học được từ các lần sản phẩm $i$ đã được mua, nên sản phẩm ít người mua có biểu diễn được ước lượng từ rất ít dữ liệu, còn sản phẩm chưa có lần mua nào thì không có biểu diễn -- vấn đề khởi đầu lạnh của sản phẩm \parencite{schein2002coldstart}; những gì đã biết về chính sản phẩm (thuộc tính, mô tả, doanh số gần đây) và về khách hàng không được dùng. Giới hạn này không thể giải quyết chỉ bằng cách làm mô hình phức tạp hơn mà phải bổ sung thông tin.

\subsection{Lựa chọn mô hình: NeuMF bổ sung đặc trưng (NeuMF-F)}
\label{subsec:sumenh_neumf}

NeuMF -- do He và cộng sự đề xuất năm 2017 \parencite{he2017ncf} -- giải quyết giới hạn thứ nhất bằng cách kết hợp thay vì thay thế MF: một nhánh tương đương MF dưới dạng mạng nơ-ron (Generalized Matrix Factorization, GMF) giữ khả năng học tương tác tuyến tính, và một nhánh Multi-Layer Perceptron (MLP) song song học thêm tương tác phi tuyến; hai nhánh được hợp nhất ở lớp dự đoán. Tuy nhiên, giai đoạn phát triển của đề tài (mục~\ref{sec:v1_history}) cho thấy NeuMF chỉ dùng mã ID không vượt được MF được tinh chỉnh tốt trên dữ liệu H\&M -- phù hợp với các tái đánh giá độc lập \parencite{rendle2020ncfmf, dacrema2019progress}.

Đề tài vì vậy xây dựng \textbf{NeuMF-F}: giữ nguyên cấu trúc hai nhánh GMF và MLP của NeuMF, nhưng biểu diễn mỗi khách hàng và sản phẩm bằng tổng của Embedding ID và một phép chiếu của đặc trưng (thuộc tính và mô tả sản phẩm, thông tin khách hàng, doanh số và thời gian), theo hướng các mô hình lai dựa trên đặc trưng như Wide \& Deep \parencite{cheng2016wide} và DeepFM \parencite{guo2017deepfm}. Phần đặc trưng đưa vào mô hình những thông tin mà mã ID không mang: sản phẩm thuộc loại gì, đang bán chạy hay không, khách hàng thuộc nhóm nào và mua gần đây hay không. Cấu trúc hai nhánh tường minh cho phép \emph{tách bạch} từng đóng góp bằng các phép so sánh trực tiếp: với NeuMF chỉ dùng ID (đóng góp của đặc trưng), với từng nhánh đứng riêng (đóng góp của việc hợp nhất), và với Late Fusion của cùng hai nhánh (hợp nhất sớm hay muộn).

\section{Mục tiêu và phạm vi nghiên cứu}
\label{sec:muctieu_phamvi}

\subsection{Mục tiêu nghiên cứu}
\label{subsec:muctieu}

\begin{itemize}
    \item Nghiên cứu cơ sở lý thuyết về Lọc cộng tác, Nhân tử hóa ma trận, Mạng nơ-ron sâu, kiến trúc NeuMF, mô hình lai dựa trên đặc trưng và hai cách hợp nhất Early Fusion, Late Fusion \parencite{he2017ncf, koren2009matrix, holinh2023multisource}.
    \item Xây dựng quy trình xử lý dữ liệu hoàn chỉnh trên bộ dữ liệu H\&M \parencite{hm2022kaggle}: lấy mẫu theo khách hàng, xây dựng ma trận tương tác phản hồi ẩn, chia dữ liệu theo thời gian, xây dựng đặc trưng sản phẩm, khách hàng và thời gian \emph{theo thời điểm} -- không dùng bất kỳ thông tin nào xảy ra sau thời điểm cần dự đoán.
    \item Thiết kế và cài đặt mô hình lai NeuMF-F trên nền tảng PyTorch \parencite{paszke2019pytorch}, cùng các mô hình đối chiếu trên cùng khung (GMF-F, MLP-F, LateFusion-F, NeuMF chỉ dùng ID) và sáu baseline (Most Popular, MostPopular-Recent, gợi ý theo nội dung, ItemKNN, UserKNN, BPR-MF), mọi mô hình được tinh chỉnh với cùng ngân sách.
    \item Đánh giá thực nghiệm định lượng theo giao thức Full Ranking trên toàn bộ tập ứng viên, bằng các độ đo NDCG@K \parencite{jarvelin2002}, Recall@K, HR@K, Precision@K; kết luận trên một mẫu khách hàng độc lập chưa từng được dùng để ra quyết định, theo kế hoạch đăng ký trước, với 5 lần chạy và kiểm định thống kê có hiệu chỉnh đa so sánh; ablation từng nhóm đặc trưng.
    \item Xây dựng hệ thống demo minh họa mô hình đã huấn luyện trong một kịch bản cửa hàng thời trang.
\end{itemize}

\subsection{Phạm vi nghiên cứu}
\label{subsec:phamvi}

\begin{itemize}
    \item \textbf{Đối tượng nghiên cứu:} Kiến trúc Neural Collaborative Filtering (GMF, MLP, NeuMF) \parencite{he2017ncf} và mở rộng có đặc trưng của nó (NeuMF-F); hai cách hợp nhất mô hình MF với mạng nơ-ron sâu -- Early Fusion (mức biểu diễn) và Late Fusion (mức điểm dự đoán) \parencite{holinh2023multisource}; các phương pháp đánh giá hệ thống gợi ý Top-K.
    \item \textbf{Dữ liệu thực nghiệm:} Bộ dữ liệu H\&M Personalized Fashion Recommendations \parencite{hm2022kaggle}: tệp giao dịch, tệp danh mục sản phẩm (thuộc tính, mô tả văn bản) và tệp thông tin khách hàng. Khảo sát quy mô được chạy trên toàn bộ 31,8 triệu giao dịch; huấn luyện và đánh giá được thực hiện trên hai mẫu khoảng 500.000 giao dịch, mỗi mẫu khoảng 21.000 khách hàng được chọn ngẫu nhiên theo hàm băm, không có khách chung, giữ trọn lịch sử 2 năm (lý do ở mục~\ref{subsec:hm_sampling}). Mỗi sản phẩm là một mẫu thiết kế (\texttt{product\_code}): mọi biến thể màu của cùng một mẫu được gộp lại. Ảnh sản phẩm không được dùng làm đầu vào mô hình.
    \item \textbf{Tín hiệu tương tác:} Phản hồi ẩn dạng nhị phân từ lịch sử mua hàng (mua hay chưa mua), kèm ngày mua.
    \item \textbf{Khởi đầu lạnh:} sản phẩm chưa có người mua trước thời điểm dự đoán (trong đó có sản phẩm vừa ra mắt) và người dùng chưa có lịch sử mua đều nằm ngoài phạm vi đánh giá định lượng (mục~\ref{subsec:split_sampling}).
\end{itemize}

\section{Phương pháp nghiên cứu}
\label{sec:phuongphap}

Để giải quyết các mục tiêu đã đặt ra (mục~\ref{subsec:muctieu}), đề tài kết hợp bốn nhóm phương pháp, được triển khai tuần tự theo trình tự trình bày trong các chương tiếp theo:

\begin{enumerate}
    \item \textbf{Phương pháp nghiên cứu lý thuyết (phân tích và tổng hợp tài liệu):} Hệ thống hóa cơ sở lý thuyết về Lọc cộng tác, Nhân tử hóa ma trận, kiến trúc NCF/NeuMF và mô hình lai dựa trên đặc trưng từ các công trình nền tảng \parencite{koren2009matrix, he2017ncf, rendle2010fm, cheng2016wide}, đồng thời khảo sát các hướng nghiên cứu hiện đại hơn để xác định vị trí đóng góp và giới hạn của hướng tiếp cận được lựa chọn (Chương~\ref{chap:lythuyet}). Các công trình được dùng làm cơ sở về phương pháp; số liệu của chúng không được dùng làm chuẩn so sánh vì khác dữ liệu và giao thức.
    \item \textbf{Phương pháp xử lý dữ liệu định lượng:} Xây dựng quy trình tiền xử lý tường minh, tái lập được và kiểm chứng được ở từng bước -- lấy mẫu theo khách hàng, gộp giao dịch, lọc $k$-core chỉ trên dữ liệu trước mốc kiểm thử, chia giai đoạn theo thời gian, xây dựng đặc trưng theo thời điểm -- với bộ kiểm thử tự động xác nhận không rò rỉ thông tin tương lai (Chương~\ref{chap:thietke_trienkhai}).
    \item \textbf{Phương pháp thực nghiệm khoa học (mô hình hóa và huấn luyện):} Cài đặt các mô hình bằng PyTorch \parencite{paszke2019pytorch}, huấn luyện bằng hàm mất mát Binary Cross-Entropy với lấy mẫu âm theo thời điểm, tối ưu bằng Adam \parencite{kingma2014adam}; siêu tham số của mọi mô hình được tinh chỉnh với cùng ngân sách, chỉ trên tập xác thực của mẫu phát triển.
    \item \textbf{Phương pháp đánh giá định lượng và kiểm chứng thống kê:} Đánh giá Full Ranking trên toàn bộ tập ứng viên; đăng ký trước kế hoạch đánh giá và khoá tập kiểm thử; kết luận trên một mẫu khách hàng độc lập; kiểm định Wilcoxon theo người dùng kèm khoảng tin cậy bootstrap và hiệu chỉnh Holm \parencite{wilcoxon1945, efron1993bootstrap, holm1979} (Chương~\ref{chap:thucnghiem}).
\end{enumerate}

Điểm xuyên suốt của phương pháp luận là ưu tiên tính tái lập và kiểm chứng được: mọi số liệu kết quả trong báo cáo đều được trích xuất tự động từ tệp kết quả do chính quy trình thực nghiệm sinh ra, không nhập tay.

\section{Ý nghĩa khoa học và thực tiễn của đề tài}
\label{sec:ynghia}

\subsection{Ý nghĩa khoa học}
\label{subsec:ynghia_khoahoc}

\begin{itemize}
    \item Cung cấp kết quả thực nghiệm có kiểm định thống kê, trên dữ liệu giao dịch thời trang thực tế và trên một mẫu khách hàng chưa từng dùng để ra quyết định, về việc mô hình lai MF + DNN có đặc trưng (NeuMF-F) có vượt các baseline được tinh chỉnh cùng ngân sách, vượt NeuMF chỉ dùng ID và vượt từng nhánh thành phần của nó hay không (kết quả ở Chương~\ref{chap:thucnghiem}).
    \item Làm rõ vị trí của NeuMF/NeuMF-F trong phân loại Early Fusion/Late Fusion (mục~\ref{subsec:early_late_fusion}) và so sánh trực tiếp hai cách hợp nhất trên cùng hai thành phần.
    \item Xây dựng một quy trình đánh giá chống thiên lệch có thể tái sử dụng cho các nghiên cứu hệ gợi ý khác: dữ liệu đã lọc ghi ra tệp kiểm bằng mã MD5, đặc trưng theo thời điểm, cùng ngân sách tinh chỉnh, mẫu kiểm định độc lập, đăng ký trước, khoá tập kiểm thử kèm nhật ký truy cập, kiểm định cặp có hiệu chỉnh đa so sánh.
\end{itemize}

\subsection{Ý nghĩa thực tiễn}
\label{subsec:ynghia_thucte}

\begin{itemize}
    \item Mô hình chỉ dùng dữ liệu mà một doanh nghiệp bán lẻ sẵn có (lịch sử giao dịch, danh mục sản phẩm, thông tin hội viên, doanh số theo ngày), và mọi đặc trưng được tính theo thời điểm, nên cách huấn luyện và chấm điểm dùng được nguyên như khi triển khai thật.
    \item Quy trình tiền xử lý tách phần đọc dữ liệu khỏi logic xử lý cốt lõi, đọc tệp theo lô (khoảng 0,5 GB RAM cho tệp 3,5 GB), nên thực nghiệm trên dữ liệu giao dịch lớn khả thi trên máy tính cá nhân và có thể tái sử dụng cho các bộ dữ liệu thương mại điện tử khác.
    \item Hệ thống demo suy diễn (FastAPI) cho phép so sánh trực quan các mô hình trên từng khách hàng của mẫu kiểm định.
\end{itemize}

\section{Bố cục của báo cáo dự án}
\label{sec:bocuc}

Báo cáo dự án được trình bày trong 5 chương chính:
\begin{itemize}
    \item \textbf{Chương 1 -- Tổng quan về đề tài:} bối cảnh, tính cấp thiết, mục tiêu, phạm vi, phương pháp nghiên cứu và ý nghĩa của đề tài.
    \item \textbf{Chương 2 -- Cơ sở lý thuyết và các công trình liên quan:} hệ thống gợi ý, Nhân tử hóa ma trận, BPR, GMF, MLP, NeuMF, Early/Late Fusion, mô hình lai dựa trên đặc trưng và khởi đầu lạnh, các chỉ số đánh giá và tổng quan các nghiên cứu liên quan.
    \item \textbf{Chương 3 -- Phương pháp nghiên cứu và thiết kế mô hình lai NeuMF-F:} dữ liệu H\&M và hai mẫu khách hàng, giai đoạn đánh giá và tập ứng viên, đặc trưng theo thời điểm, kiến trúc NeuMF-F và các mô hình đối chiếu, lấy mẫu âm, huấn luyện, quy trình tinh chỉnh và đánh giá không thiên lệch, thiết kế hệ thống demo.
    \item \textbf{Chương 4 -- Thực nghiệm và đánh giá kết quả:} kết quả tinh chỉnh trên mẫu phát triển, kết quả đánh giá cuối trên mẫu kiểm định, kiểm định thống kê, ablation, độ phủ, lịch sử phát triển và thảo luận.
    \item \textbf{Chương 5 -- Kết luận và hướng phát triển:} tổng kết kết quả, hạn chế và hướng mở rộng.
\end{itemize}
